% Source: Wald GR, physical PDF pp. 39--41
%         (printed pp. 26--28).
% Coverage: Chapter 2 problems 1--8 and end material.
% Figures/tables/footnotes: no figures, tables, or footnotes.
% Status: source-reviewed and compile-clean.
% Uncertainties: none.

\chapterproblems

\begin{enumerate}[label=\arabic*., leftmargin=*]
  \item
  \begin{enumerate}[label=\alph*)]
    \item Show that the overlap functions
      \(f_i^\pm\circ(f_j^\pm)^{-1}\) are \(C^\infty\), thus completing the
      demonstration given in Section~\ref{sec:2.1} that \(S^2\) is a manifold.

    \item Show by explicit construction that two coordinate systems (as opposed to
      the six used in the text) suffice to cover \(S^2\).  (It is impossible to
      cover \(S^2\) with a single chart, as follows from the fact that \(S^2\) is
      compact, but every open subset of \(\mathbb{R}^2\) is noncompact; see
      appendix A.)
  \end{enumerate}

  \item Prove that any smooth function \(F:\mathbb{R}^n\to\mathbb{R}\) can be
  written in the form Eq.~\eqref{eq:2.2.2}.  (Hint: For \(n=1\), use the identity
  \[
    F(x)-F(a)=(x-a)\int_0^1 F'[t(x-a)+a]\,dt ;
  \]
  then prove it for general \(n\) by induction.)

  \item
  \begin{enumerate}[label=\alph*)]
    \item Verify that the commutator, defined by Eq.~\eqref{eq:2.2.14}, satisfies
      the linearity and Leibniz properties, and hence defines a vector field.

    \item Let \(X,Y,Z\) be smooth vector fields on a manifold \(M\).  Verify
      that their commutator satisfies the Jacobi identity:
      \[
        [[X,Y],Z]+[[Y,Z],X]+[[Z,X],Y]=0 .
      \]

    \item Let \(Y_1,\ldots,Y_n\) be smooth vector fields on an \(n\)-dimensional
      manifold \(M\) such that at each \(p\in M\) they form a basis of the tangent
      space \(T_pM\).  Then, at each point, we may expand each commutator
      \([Y_\alpha,Y_\beta]\) in this basis, thereby defining the functions
      \(C^\gamma{}_{\alpha\beta}=-C^\gamma{}_{\beta\alpha}\) by
      \[
        [Y_\alpha,Y_\beta]=\sum_\gamma C^\gamma{}_{\alpha\beta}Y_\gamma .
      \]
      Use the Jacobi identity to derive an equation satisfied by
      \(C^\gamma{}_{\alpha\beta}\).  (This equation is a useful algebraic relation
      if the \(C^\gamma{}_{\alpha\beta}\) are constants, as will be the case if
      \(Y_1,\ldots,Y_n\) are left [or right] invariant vector fields on a Lie group
      [see Section~7.2].)
  \end{enumerate}

  \item
  \begin{enumerate}[label=\alph*)]
    \item Show that in any coordinate basis, the components of the commutator of
      two vector fields \(v\) and \(w\) are given by
      \[
        [v,w]^\mu=\sum_\nu\left(
        v^\nu\frac{\partial w^\mu}{\partial x^\nu}
        -w^\nu\frac{\partial v^\mu}{\partial x^\nu}\right) .
      \]

    \item Let \(Y_1,\ldots,Y_n\) be as in problem 3(c).  Let
      \(f^{1},\ldots,f^{n}\) be the dual basis.  Show that the components
      \((f^{\gamma})_\mu\) of \(f^{\gamma}\) in any coordinate basis satisfy
      \[
        \frac{\partial(f^{\gamma})_\mu}{\partial x^\nu}
        -\frac{\partial(f^{\gamma})_\nu}{\partial x^\mu}
        =\sum_{\alpha,\beta}C^\gamma{}_{\alpha\beta}
        (f^{\alpha})_\mu(f^{\beta})_\nu .
      \]
      (Hint: Contract both sides with \((Y_\alpha)^\mu(Y_\beta)^\nu\).)
  \end{enumerate}

  \item Let \(Y_1,\ldots,Y_n\) be smooth vector fields on an \(n\)-dimensional
  manifold \(M\) which form a basis of \(T_pM\) at each \(p\in M\).  Suppose
  \([Y_\alpha,Y_\beta]=0\) for all \(\alpha,\beta\).  Prove that in a neighborhood
  of each \(p\in M\) there exist coordinates \(y^1,\ldots,y^n\) such that
  \(Y_1,\ldots,Y_n\) are the coordinate vector fields,
  \(Y_\mu=\partial/\partial y^\mu\).  (Hint: In an open ball of \(\mathbb{R}^n\),
  the equations \(\partial f/\partial x^\mu=F_\mu\) with
  \(\mu=1,\ldots,n\) for the unknown function \(f\) have a solution if and only
  if \(\partial F_\mu/\partial x^\nu=\partial F_\nu/\partial x^\mu\).  [See the
  end of Section B.1 of appendix B for a statement of generalizations of this
  result.]  Use this fact together with the results of problem 4(b) to obtain the new
  coordinates.)

  \item
  \begin{enumerate}[label=\alph*)]
    \item Verify that the covectors \(\{f^{\mu}\}\) defined by
      Eq.~\eqref{eq:2.3.1} constitute a basis of \(V^\vee\).

    \item Let \(e_1,\ldots,e_n\) be a basis of the vector space \(V\), and let
      \(f^{1},\ldots,f^{n}\) be its dual basis.  Let \(w\in V\) and let
      \(\omega\in V^\vee\).  Show that
      \[
        w=\sum_\alpha f^{\alpha}(w)e_\alpha ,
        \qquad
        \omega=\sum_\alpha\omega(e_\alpha)f^{\alpha} .
      \]

    \item Prove that the operation of contraction, Eq.~\eqref{eq:2.3.2}, is
      independent of the choice of basis.
  \end{enumerate}

  \item Let \(V\) be an \(n\)-dimensional vector space and let \(g\) be a metric
  on \(V\).
  \begin{enumerate}[label=\alph*)]
    \item Show that one always can find an orthonormal basis \(e_1,\ldots,e_n\) of
      \(V\), i.e., a basis such that \(g(e_\alpha,e_\beta)=\pm\delta_{\alpha\beta}\).
      (Hint: Use induction.)

    \item Show that the signature of \(g\) is independent of the choice of
      orthonormal basis.
  \end{enumerate}

  \item
  \begin{enumerate}[label=\alph*)]
    \item The metric of flat, three-dimensional Euclidean space is
      \[
        ds^2=dx^2+dy^2+dz^2 .
      \]
      Show that the metric components \(g_{\mu\nu}\) in spherical polar
      coordinates \(r,\theta,\phi\) defined by
      \[
        r=(x^2+y^2+z^2)^{1/2},
        \qquad \cos\theta=z/r,
        \qquad \tan\phi=y/x
      \]
      is given by
      \[
        ds^2=dr^2+r^2d\theta^2+r^2\sin^2\theta\,d\phi^2 .
      \]

    \item The spacetime metric of special relativity is
      \[
        ds^2=-dt^2+dx^2+dy^2+dz^2 .
      \]
      Find the components, \(g_{\mu\nu}\) and \(g^{\mu\nu}\), of the metric and
      inverse metric in ``rotating coordinates,'' defined by
      \[
        \begin{aligned}
          t'&=t,\\
          x'&=(x^2+y^2)^{1/2}\cos(\phi-\omega t),\\
          y'&=(x^2+y^2)^{1/2}\sin(\phi-\omega t),\\
          z'&=z,
        \end{aligned}
      \]
      where \(\tan\phi=y/x\).
  \end{enumerate}
\end{enumerate}
